\documentclass[10pt]{article}
\usepackage[T1]{fontenc}
\usepackage[utf8]{inputenc}
\usepackage[french]{babel}
\usepackage{amsmath,amssymb,mathrsfs,mathpazo}
\usepackage{geometry}
\geometry{a4paper,margin=2cm}
\usepackage{graphicx,array,tabularx,tabularray,colortbl,multirow,multicol,booktabs,enumitem,pifont,ulem,fancybox,diagbox,cancel}
\usepackage{pstricks,pstricks-add,pst-plot,pst-tree,pst-node,pst-text,pst-func,pst-math,pst-bspline,pst-3dplot,pst-eucl,pst-solides3d}
\let\pscalculate\undefined
\usepackage{pst-bezier}
\usepackage{tikz,pgfplots,tkz-tab}
\usetikzlibrary{arrows,arrows.meta,patterns,calc,positioning,trees,decorations.markings,shapes.geometric,shapes.misc,shapes.arrows,intersections,decorations.pathmorphing,decorations.pathreplacing,backgrounds,shadows,babel}
\pgfplotsset{compat=1.18}
\usepgfplotslibrary{fillbetween}
\usepackage[np]{numprint}
\usepackage{listings,makecell,pas-tableur,twoopt}
\let\barre\undefined
\usepackage[pstricks]{tablvar}
\definecolor{darkgreen}{RGB}{0,100,0}
\providecommand{\red}{\color{red}}
\providecommand{\blue}{\color{blue}}
\providecommand{\green}{\color{green}}
\providecommand{\black}{\color{black}}
\providecommand{\white}{\color{white}}
\providecommand{\gray}{\color{gray}}
\providecommand{\cyan}{\color{cyan}}
\providecommand{\magenta}{\color{magenta}}
\providecommand{\yellow}{\color{yellow}}
\providecommand{\violet}{\color{violet}}
\providecommand{\orange}{\color{orange}}
\providecommand{\euro}{\texteuro}
\providecommand{\EUR}{\texteuro}
\providecommand{\ieme}{\textsuperscript{e}}
\providecommand{\degres}{\ensuremath{^\circ}}
\providecommand{\degree}{\ensuremath{^\circ}}
\providecommand{\vv}[1]{\overrightarrow{#1}}
\providecommand{\Coord}[1]{#1}
\providecommand{\decosix}{}
\providecommand{\starredbullet}{\ensuremath{\bullet}}
\pagestyle{empty}
\setlength{\parindent}{0pt}
\usepackage{adjustbox}\DeclareUnicodeCharacter{2236}{:}\DeclareUnicodeCharacter{2190}{\ensuremath{\leftarrow}}
\providecommand{\R}{\mathbb{R}}
\providecommand{\N}{\mathbb{N}}
\providecommand{\Z}{\mathbb{Z}}
\providecommand{\Q}{\mathbb{Q}}
\providecommand{\D}{\mathbb{D}}
\providecommand{\C}{\mathbb{C}}
\providecommand{\e}{{\mathrm{e}}}
\providecommand{\vect}[1]{\overrightarrow{#1}}
\providecommand{\vectt}[1]{\overrightarrow{\mathrm{#1}}}
\providecommand{\vvt}[1]{\overrightarrow{\mathrm{#1}}}
\providecommand{\np}[2][]{{#2~#1}}
\providecommand{\vv}[1]{{\overrightarrow{#1}}}
\renewcommand{\d}{{\mathrm{d}}}
\providecommand{\degres}{{^\circ}}
\providecommand{\degre}{{^\circ}}
\providecommand{\blue}{\color{blue}}
\providecommand{\red}{\color{red}}
\providecommand{\enskip}{\quad}
\providecommand{\llbracket}{\lbrack\!\lbrack}
\providecommand{\rrbracket}{\rbrack\!\rbrack}
\providecommand{\euro}{{\text{€}}}
\providecommand{\no}{{\text{n°}}}
\providecommand{\interval}[2]{{[#1;#2]}}
\providecommand{\intervFO}[2]{{[#1;#2[}}
\providecommand{\dotfill}{\ldots}

\providecommand{\ds}{\displaystyle}
\providecommand{\textding}[1]{#1}
\providecommand{\Oij}{\(\mathrm{O};\vec{\imath},\vec{\jmath}\)}
\providecommand{\Oijk}{\(\mathrm{O};\vec{\imath},\vec{\jmath},\vec{k}\)}
\providecommand{\Ouv}{\(\mathrm{O};\vec{u},\vec{v}\)}
\providecommand{\intervFF}[2]{\ensuremath{[#1;#2]}}
\renewcommand{\euro}{\ifmmode\text{\texteuro}\else\texteuro\fi}
\DeclareUnicodeCharacter{00B0}{\ifmmode^{\circ}\else\textdegree\fi}
\DeclareUnicodeCharacter{20AC}{\ifmmode\text{\texteuro}\else\texteuro\fi}
% Énoncés : Éducation nationale. Transcription et corrigés : APMEP, Denis Vergès et contributeurs. Découpage et classement : CoopMaths, Sébastien Lozano, Éric Elter et contributeurs. Adaptation de la présentation : Maths Explorer.
% CC BY-SA : https://coopmaths.fr/www/about/
% https://coopmaths.fr/alea/static/bac/2022/tex/bac_2022_10_sujet2_caledonie_4_cor.tex
\begin{document}
\noindent\textit{APMEP / CoopMaths. CC BY-SA. Présentation adaptée par Maths Explorer.}\par\medskip
\textbf{Principaux domaines abordés :} suites, fonctions, primitives

\medskip









\begin{enumerate}
\item 

\[u_n = \dfrac{(- 1)^n}{n + 1}.\]








On a $- 1 \leqslant (- 1)^n \leqslant 1$, donc puisque $n+1 \geqslant 1 > 0$ et que $\dfrac{1}{n  1} > 0$, on en déduit que pour tout naturel :

\[- \dfrac{1}{n+1} \leqslant u_n \leqslant \dfrac{1}{n+1}\]

La suite $\left(u_n\right)$ encadrée par deux suites ayant pour limite 0 a pour limite 0.
\begin{center} \decosix \decosix \decosix  \end{center}
\end{enumerate}

Dans les questions 2 et 3, on considère deux suites $\left(v_n\right)$  et $\left(w_n\right)$ vérifiant la relation : 

\[w_n = \e^{- 2v_n} + 2.\]










On a $w_0 = \e^{- 2v_0} + 2 = \e^{- 2\ln (a)} + 2 = \dfrac{1}{\e^{2\ln(a)}} + 2=  \dfrac{1}{\e^{\ln \left(a^2\right)}} + 2 = \dfrac{1}{a^2} + 2$.

\begin{enumerate}[resume]
\item On sait que la suite $\left(v_n\right)$ est croissante. On peut affirmer que la suite $\left(w_n\right)$ est :







On a successivement :

$v_n \leqslant v_{n+1} \Rightarrow 2v_n \leqslant 2v_{n+1} \Rightarrow - 2v_{n+1} \leqslant - 2v_n\Rightarrow \e^{- 2v_{n+1}} \leqslant \e^{- 2v_n}$ par croissance de la fonction exponentielle et enfin $\e^{- 2v_{n+1}} + 2 \leqslant \e^{- 2v_n} + 2$, soit $w_{n+1} \leqslant w_n$ : la suite $\left(w_n\right)$ est décroissante.

D'autre part on sait que quel que soit le réel$\alpha, \: \e^{\alpha} > 0$ donc  la suite est minorée par 2.
\item 

\[a_0 = 2 \:\:\text{et, pour tout entier naturel}\: n, \quad a_{n+1} = \dfrac13a_n + \dfrac83.\]








Soit la suite $\left(b_n\right)$ définie pour tout naturel $n$ par $b_n = a_n + \alpha$.

Alors $b_{n+1} = a_{n+1} + \alpha = \dfrac13a_n + \dfrac83 + \alpha = \dfrac13\left(b_n - \alpha \right) + \dfrac83 + \alpha = \dfrac13b_n - \dfrac{\alpha}{3} + \dfrac83 + \alpha = \dfrac13b_n + \dfrac{8 + 2\alpha}{3}$.

Prenons $\alpha = - 4$, alors $b_{n+1} = \dfrac13b_n$ : cette égalité montre que la suite $\left(b_n\right)$ est une suite géométrique de raison $\dfrac13$ et de premier terme $b_0 = a_0 + \alpha = 2 - 4 = - 2$.

On sait qu'alors que, quel que soit le naturel $n$, \: $b_n = b_0 \times \left(\dfrac{1}{3}\right)^n = - 2 \times \left(\dfrac{1}{3}\right)^n$.

Finalement $b_n = a_n - 4 \iff a_n = 4 + b_n = 4 - 2\times \left(\dfrac{1}{3}\right)^n$.
\item On considère une suite $\left(b_n\right)$ telle que, pour tout entier naturel $n$, on a :

\[b_{n+1} = b_n + \ln \left(\dfrac{2}{\left(b_n \right)^2 + 3}\right).\]









On a quel que soit le naturel $n$, \: $b_{n+1} - b_n = \ln \left(\dfrac{2}{\left(b_n \right)^2 + 3}\right)$.

Or quel que soit le réel $b_n$, \: $b_n^2 \geqslant 0$, donc $\left(b_n\right)^2 + 3 \geqslant 3 > 2$, donc $\dfrac{2}{\left(b_n \right)^2 + 3} < 1$ et enfin $\ln \left(\dfrac{2}{\left(b_n \right)^2 + 3}\right) < 0$.

Conclusion : $b_{n+1} - b_n < 0$ montre que la suite $\left(b_n\right)$ est décroissante.

\item 

\[g(x) = \dfrac{\e^x}{x}.\]













$\bullet~~$Au voisinage de zéro : $\displaystyle\lim_{x \to 0} \e^x = 1$ et $\displaystyle\lim_{x \to 0}x  = 0$, donc $\displaystyle\lim_{x \to 0} g(x) = +\infty$ : géométriquement l'axe des ordonnées est asymptote verticale à $\mathcal{C}_g$.

$\bullet~~$Au voisinage de plus l'infini : on sait que $\displaystyle\lim_{x \to + \infty} g(x) = + \infty$, donc pas d'asymptote.

\item 

\[f(x) = x\e^{x^2+1}.\]









En posant $u(x) = x^2 + 1$, on a $u'(x) = 2x$ et l'on peut écrire $f(x) = \dfrac12 \times 2x \times \e^{x^2+1} = \dfrac12 u'\e^u$ : on reconnait la dérivée de $\dfrac12\e^u = \dfrac12\e^{x^2 + 1}$.

Donc $F(x) = \dfrac12\e^{x^2+1}$.
\end{enumerate}

\end{document}
